\newpage
\section{Conclusion}
\label{sec:conclusion}

This thesis set out from a deceptively simple observation---that a Kubernetes pod is a unit
of scheduling whereas a Redis shard is a unit of data ownership, and that the two coincide
only for as long as the cluster's topology remains undisturbed---and asked what happens to a
sharded, stateful workload when that topology is disturbed in the most consequential way, by
removing a slot-owning master during a planned scale-in. Through a deliberately minimal and
fully controlled experiment, it has shown that the answer is governed not by how a scaling
decision is timed, nor by whether redundancy is present, nor by the load the cluster happens
to be under, but by a single architectural property of the mechanism that carries the removal
out: whether or not that mechanism is aware of the cluster's slot topology.

\subsection{Revisiting the Research Question}

The central question---whether a planned, one-master scale-in of a sharded Redis Cluster
preserves all \num{16384} hash slots together with every stored key, and to what extent the
answer depends on the scaling mechanism---admits a clear and sharply divided answer. When the
removal was performed by the native StatefulSet contraction, which is precisely the action a
Horizontal Pod Autoscaler issues on scale-down, the cluster suffered an immediate and
irreversible \textit{data cliff}: the departing master's slots, neither drained nor
resharded, were abandoned the instant its pod terminated, slot coverage fell below
\num{16384}, the cluster state transitioned to \texttt{fail}, roughly a quarter of the
keyspace became permanently unreachable, and---under the default
\texttt{cluster-require-full-coverage} policy---every read was refused with a
\texttt{CLUSTERDOWN} error. When the same transition was requested of the Opstree Redis
Operator, the outcome was the exact opposite: the operator drained the departing master,
migrating its slots to the survivors before removing the pod, so that full coverage was held
throughout and not a single key was lost. The hypothesis advanced in
Section~\ref{sec:hypothesis}---that topology-aware operation is a precondition for data-safe
scale-in---is therefore confirmed.

Two features of this result give it particular force. First, the loss in the native case
occurred under perfectly quiescent conditions, with no client traffic of any kind, which
establishes the cliff as a \textit{structural} consequence of the mechanism rather than an
artefact of write pressure, contention, or timing. Second, the divergence between total
preservation and catastrophic loss was produced while every other variable---the hardware,
the dataset, the resource budget, and the absence of load---was held identical, which
isolates topology-awareness as the sole cause. Redundancy is no substitute: a replica would
have delayed the cliff by one promotion but could not have prevented it, because a planned
scale-in fundamentally requires the migration of slots that no native primitive performs.

\subsection{Contributions}

The principal contributions of this thesis may be stated concisely. It demonstrates, in
isolation and with a deterministic and reproducible method, that the native, topology-blind
scaling mechanism inflicts structural and irreversible data loss on the scale-in of a sharded
Redis Cluster, and that a topology-aware operator preserves the keyspace in full under
identical conditions (\cite{redis:scaling}; \cite{opstree:redisoperator}). In doing so it
establishes that the determining factor for data safety is the scaling \textit{mechanism}
rather than the scaling \textit{trigger}, thereby reframing a question that the existing
literature has approached almost exclusively through the lens of reactive-versus-proactive
timing. Methodologically, it contributes a minimal experimental design---direct dataset
population, a single quiescent one-master transition, and measurement through the cluster's
own slot and key telemetry---that strips away the confounds of a fuller pipeline and renders
the mechanism the only variable. Finally, and importantly for the credibility of the result,
it characterises the proposed solution honestly rather than uncritically: the operator is
shown to be \textit{necessary but not, on its own, sufficient} for fully hands-off elastic
operation, because its lack of self-healing for an interrupted reshard can leave a slot in an
open, mid-migration state that a human operator must settle by hand with
\texttt{redis-cli -{}-cluster fix}.

\subsection{Reflection on Scope and Methodology}

The strength of the foregoing result rests on the deliberate narrowness of the experiment
that produced it, and it is worth reflecting on that scoping explicitly. The study was
confined to the data-safety of a single scale-in under quiescent conditions, and this
confinement was not an evasion of difficulty but a direct response to one. An earlier,
fuller pipeline---comprising an external load generator, an application entrypoint, and an
event-driven autoscaling trigger---was found to confound the very measurement it was meant to
support: under a sustained, write-heavy workload, a continuous stream of writes addressed to
a slot that was mid-migration was the dominant cause of an interrupted reshard, and an
interruption so caused would have been indistinguishable from a genuine defect of the scaling
mechanism. Removing the load, the entrypoint, and the trigger, and writing the dataset
directly, eliminated that confound and made the outcome a deterministic property of the
mechanism alone. The load-induced interruption is not thereby dismissed; it is itself a
finding, namely that write pressure rather than the mechanism is the governing factor in
reshard interruption, and it is what justifies the controlled isolation.

It follows that certain elements present in the thesis's initial conception were, in the
final analysis, deliberately reframed as out of scope rather than evaluated. The event-driven
trigger layer in particular---the question of whether scaling decisions are taken reactively
or proactively, and the role of application-level metrics in that decision---was set aside as
orthogonal to the data-safety question, on the reasoning that no quality of trigger can
rescue an unsafe removal mechanism and none is needed to vindicate a safe one. What remains,
and what the thesis ultimately defends, is a tightly bounded but firmly established claim
about the scaling mechanism itself. The remaining limitations follow from the same boundary
and are enumerated among the threats to validity in Section~\ref{sec:results}: the operator's
robustness under concurrent load is not characterised here, the native cluster was formed
without replicas so as to expose the cliff in its starkest form, and the comparison concerns
a single one-master transition rather than a repeated or multi-master contraction.

\subsection{Future Work}

These boundaries map directly onto the most promising directions for further work. The most
immediate is the reintroduction of live load under controlled conditions, in order to
characterise how frequently the operator's reshard is interrupted under realistic write
pressure and how that frequency scales with throughput---turning the confound identified
above into an object of study in its own right. A natural second direction is to layer an
event-driven trigger, such as the Kubernetes Event-Driven Autoscaling project, atop the
topology-aware mechanism, so as to combine a proactive scaling decision with a data-safe
scaling action and to evaluate the pairing against the reactive baseline. A third, and
perhaps the most consequential for practice, is to close the self-healing gap that renders
the operator necessary but not sufficient: automating the detection and repair of an
interrupted reshard would move the operator from data-safe-but-supervised toward genuinely
hands-off operation. Beyond these, the comparison invites extension to repeated and
multi-master contractions, to clusters that enable replication and persistence, to
substantially larger datasets, and to other operators and other sharded stores, so as to test
the generality of the conclusion drawn here.

\subsection{Concluding Remarks}

The title of this thesis asserts that pods are not shards, and the experiment has given that
assertion an empirical edge. Treating the removal of a stateful, slot-owning node as though
it were the removal of an interchangeable, stateless replica does not merely degrade
performance; it destroys data, structurally and irreversibly, and it does so the moment the
pod is deleted. The remedy is not a faster or more anticipatory autoscaler but an autoscaling
\textit{mechanism} that understands what it is removing---one that drains before it deletes.
Such topology-aware operation, as embodied here by a Kubernetes operator, is therefore a
necessary foundation for the safe elastic operation of sharded, stateful workloads; the
remaining distance to fully autonomous operation is measured precisely by the operational toil
that the operator, as it stands today, still leaves to a human hand.
